% theme: flowering_plants
\MajorQuestion{植物のなかま（1）}
\SetRowSpaceQanda

\Instruction{次の問いに答えなさい。}
\QQ{花の中心にあり、先端と根元に特徴的なつくりをもつ部分を\NamedBlank{pistil}という。}{めしべ}
\QQ{\RefBlank{pistil}の根元にある、ふくらんだ部分を\NamedBlank{ovary}という。}{子房}
\QQ{\RefBlank{pistil}の先端にある部分を\NamedBlank{stigma}という。}{柱頭}
\QQ{\RefBlank{pistil}のまわりにあり、先端で花粉をつくる部分を\NamedBlank{stamen}という。}{おしべ}
\QQ{\RefBlank{ovary}の中にある小さな粒を\NamedBlank{ovule}という。}{胚珠}

\Instruction{\strut}
\QQ{\RefBlank{stamen}の外側にあり、目立つ色や形をした部分を\NamedBlank{petal}という。}{花弁}
\QQ{\RefBlank{stamen}の先端にある、袋のような部分を\NamedBlank{anther}という。}{やく}
\QQ{\RefBlank{petal}の外側にあり、花を保護する部分を何というか。}{がく}
\QQ{\RefBlank{anther}の中でつくられる粉を\NamedBlank{pollen}という。}{花粉}
\QQ{\RefBlank{pollen}が\RefBlank{stigma}につくことを\NamedBlank{pollination}という。}{受粉}

\Instruction{\strut}
\QQ{\RefBlank{pollination}のあと、\RefBlank{ovary}が成長したものを何というか。}{果実}
\QQ{\RefBlank{pollination}のあと、\RefBlank{ovule}が成長したものを\NamedBlank{seed}という。}{種子}
\QQ{花を咲かせ、\RefBlank{seed}をつくってふえる植物を\NamedBlank{seedplant}という。}{種子植物}
\QQ{\RefBlank{seedplant}のうち、\RefBlank{ovule}が\RefBlank{ovary}の中にあるものを\NamedBlank{angiosperm}という。}{被子植物}
\QQ{\RefBlank{angiosperm}のうち、発芽したときの子葉が1枚のものを\NamedBlank{monocot}という。}{単子葉類}

\Instruction{\strut}
\QQ{\RefBlank{monocot}に多い葉脈を何というか。}{平行脈}
\QQ{\RefBlank{angiosperm}のうち、発芽したときの子葉が2枚のものを\NamedBlank{dicot}という。}{双子葉類}
\QQ{\RefBlank{monocot}に多い、1か所から多数の細い根が出るつくりを何というか。}{ひげ根}
\QQ{\RefBlank{dicot}に多い葉脈を何というか。}{網状脈}
\QQ{\RefBlank{dicot}の根は、1本の太い根と枝分かれした細い根からなる。この2つを何というか。}{主根と側根}
